% !PRINTED 365
\weberpage{365}
%<<CONT
groups also by $a'B$, where $a'$ is contained in $A'$. If then $a_1'B, a_2'B$ are \gesperrt{different} subsidiary groups to $B$ in $A$, then $a_1'B', a_2'B'$ are \gesperrt{different} subsidiary groups to $B'$ in $A'$, and conversely; for $a_1'a_2'^{-1}$ is contained in $B'$ then and only then, when it is at the same time contained in $B$, and from this follows formula (9).

We now apply these theorems to the expressions (2) and (3) for $h$ and $h'$. By (4) we have
\weq{10}{(\overline{O}, \overline{E}O') = (\overline{O}, \overline{E}O)(\overline{E}O, \overline{E}O'),}
hence
\weq{11}{h' = h(\overline{E}O, \overline{E}O').}

We now denote, as already above, by $E$ the group of the numerical units in $O$. Then $\overline{E}E = \overline{E}$, since the numerical units are contained among the functional ones, and every \gesperrt{number} that is contained in $\overline{E}O'$ is contained in $EO'$. Consequently $EO'$ is the intersection of $O$ and $\overline{E}O'$. If we therefore apply (7), putting $O, EO', \overline{E}$ in place of $A, B, C$, then $B$ is the intersection of $A$ and $BC$, and we obtain:
\weq{12}{(\overline{E}O, \overline{E}O') = (\overline{E}O, \overline{E}EO') = (O, EO').}

Further, by (4):
\weq{13}{(O, EO')(EO', O') = (O, O'),}
and if we denote by $E'$ the group of the numerical unit in $O'$ % [sic] German source prints the singular Einheit., that is, the intersection of $E$ with $O'$, so that $E'O' = O'$:
\[
(EO', O') = (EO', E'O') = (E, E')
\]
(by 2), for $E'$ is the intersection of $E'O'$ and $E$.

Thus by (13) we have
\[
(O, EO')(E, E') = (O, O')
\]
and by (12)
\[
(\overline{E}O, \overline{E}O')(EE') = (O, O'), % [sic] German source prints no comma between E and E'.
\]
hence finally by (11):
\weq{14}{(E, E')h' = (O, O')h.}

15. \gesperrt{Theorem. Formula (14) holds unchanged if $O'$ denotes not precisely the order group, but any number group of finite index $(O, O')$ contained in $O$, provided we call two ideals $\mathfrak{a}, \mathfrak{a}'$ equivalent according to $O'$ if their quotient $\mathfrak{a}'/\mathfrak{a} = \eta'$ is a number in $O'$.}
